\documentclass[../main.tex]{subfiles}

\begin{document}

\ifSubfilesClassLoaded{

    \tableofcontents
    
}{}

\chapter{ Countability }

\begin{definition}{Second-Countable Space}{}
     topological space $X$ is called second-countable if it has a countable basis for its topology. 
\end{definition}

\begin{definition}{Lindel\"of Space}{}
 A topological space $X$ is called a Lindelöf space if every open cover of $X$ has a countable subcover.
\end{definition}

\begin{theorem}{Linde\"of Theorem}{}
    Every second-countable space is Linde\"of.
\end{theorem}
\begin{proof}

    Let \(  X  \) be a second-coutable space, \(  \mathcal{B}= \left\{ B_{n} \right\}_{n = 1}^{\infty}  \) is a countable basis of \(  X  \). Let \(  \mathcal{U}= \left\{ U_{\alpha } \right\}_{\alpha \in A}  \) be any open cover of \(  X  \).     
  Let \[
  \mathcal{B}^{\prime} \coloneqq \left\{ B_{n}\in \mathcal{B}: \exists U_{\alpha }\in \mathcal{U}, s.t B_{n}\subseteq U_{\alpha } \right\}
  \] For each \(  x \in X  \),  \(  \exists U_{\alpha }\in \mathcal{U}  \) such that \(  x \in U_{\alpha }  \). Since \(  \mathcal{B}  \) is a basis, there \(  \exists   \) \(  B_{n}\in \mathcal{B}  \)       such that \(  x \in B_{n}\subseteq U_{\alpha }  \). \(  B_{n}\in \mathcal{B}^{\prime}   \) by definition. Thus \(  \mathcal{B}^{\prime}   \) covers \(  X  \). For each \(  B_{n_{k}}\in \mathcal{B}^{\prime}   \)     , we choose an open set \(  U_{k}  \) in \(  \mathcal{U}  \), such that \(  B_{n_{k}}\subseteq U_{k}  \). Then \(  \left\{ U_{k} \right\}_{k= 1}^{\infty}  \) is a subcover of \(  \mathcal{U}  \).     
\end{proof}
\end{document}